% Plan: development/outline.md, section 6, item B.5 (hvycrash, ex6_2_5,
% ex6_2_7, etamac, pricing050, pindyck).
% Round-1 revision (adjudication G5): one certificate box per result, the
% verification record merged into the boxes; Gibbs second-order bound stated
% with proof (G5-03).
\subsection{The six small instances}\label{app:small}

This subsection gives the models, the full proofs and the computations behind
\cref{prop:hvycrash-identity,thm:ex62-bounds,thm:pricing-bound,thm:etamac-bound,thm:pindyck-bound,cor:pindyck-unique}, which \cref{sec:other} states.
At the level of $10^{-14}$ and below, the stored decimals differ from the source models (the \inst{ex6_2_7} constants, the \inst{etamac} exponents, the \inst{pricing050} coefficients $g_{ij}$ and the \inst{pindyck} constant $k$; each is given below), so our enclosures need not contain the optimal values of the source models.
For each instance at least two separately written OSIL readers agree on the model.

% ---------------------------------------------------------------------------
\subsubsection{\inst{hvycrash}}\label{app:small-hvycrash}

\paragraph{Model.}
\Cref{sec:other-hvycrash} displays the rows and the objective; in the OSIL file, for $k=1,\dots,50$,
$\theta_k=x_k$ and $\theta_0=x_{101}$ (all in $[0,6.2831854]$), $c_k=x_{50+k}\in[0.08,0.417]$, and the free variables are $r_k=x_{101+k}$ and $s_k=x_{202-k}$;
the rows $\text{acc}_k$, $\text{alg}_k$, $\text{dyn}_k$ are $e_{3k-2}$, $e_{3k-1}$, $e_{3k}$, and the objective is $x_{152}=s_{50}$.
A check verified, symbolically and in exact rationals, the two identities used in the proof of \cref{prop:hvycrash-identity},
\[
\text{acc}_k=s_{k-1}-s_k-h-h\,r_k\,\text{alg}_k,\qquad
\text{dyn}_k=0.1(\theta_{k-1}-\theta_k)+\frac{h\,(A(c_k)+1)}{r_k^2}+\frac{h\,\text{alg}_k}{r_k},
\]
for all 50 stages, where each row name stands for its left side.

\paragraph{Numbers in the existence proof.}
With $c_k=0.08$, $A(0.08)=1000/149$, $D(0.08)=24149771/1250000000$ and
$\kappa=10h(A(0.08)+1)D(0.08)=81381972927/12500000000000<0.006511$.
The Taylor series of the cosine at $2.6$ alternates, and its terms decrease in absolute value from the second term on;
so the partial sum through $x^{12}$, which ends with a positive term, is an upper bound, and in exact rationals it gives $\cos2.6\le-0.8568815<-0.85$.
Then $50\kappa/0.85<0.3830$, so every angle of the backward recursion lies in $[3-0.3830,\,3]\subset[2.6,3]$;
the recursion gives $\theta_0=2.6565078758\ldots$ and $\max_k\cos\theta_k<-0.88$.
Hence every angle lies in $[0,6.2831854]$, every $c_k=0.08$ is at its lower bound, and $r_k$ and $s_k$ are free, so the witness satisfies every bound.
The witness is defined by the recursion, not by a decimal vector.

\begin{certbox}[Certificate for \cref{prop:hvycrash-identity}]
\certitem{Inputs}{OSIL file (SHA-256 \texttt{27a0a05a}).}
\certitem{Computation}{none needed; checks: identities in \arith{E}; all 150 rows and the bounds at two witnesses, 60-digit \arith{I}.}
\certitem{Data reading}{exact; row structure, including the signs of the linear parts, matched by three codes with two separately written readers.}
\certitem{Implementations}{written proof; two codes checked the first witness, and the second code proved a second witness ($c_k=0.417$, $\theta_{50}=4.2$, all $\theta_k\in[2.1653,4.2]$) in the same way.}
\certitem{Evidence level}{\evid{hand}.}
\certitem{Replay}{Tier~1; under 1\,s.}
\end{certbox}

\begin{remark}[Tolerance-feasible values]\label{rem:hvycrash-tolerance}
Under a positive tolerance any stage can be dropped as in \cref{sec:other-hvycrash}, by letting $r_k\to\infty$: then $\text{alg}_k$ and $\text{dyn}_k$ hold up to small residuals and the increment of $s$ vanishes, so tolerance-feasible points can reach values near $-0.2185+jh$ for $j=1,\dots,50$.
The COCONUT record lists the value $-0.1573199996$ with ``infeas = 14.00'' \citep{coconut2026-globallib-gams-coconut-and-minlplib}, close to $-0.2185+14h$; we did not reconstruct that point.

The identity also limits how far below $-0.2185$ tolerance-feasible points can lie.
Write $q_k=\cos\theta_k/(D(c_k)r_k^2)$, so that $\text{acc}_k$ adds $h\,q_k$ to $s$.
If every row may be violated by $\tau>0$, then $\text{alg}_k=-(1+q_k)/r_k$ gives $|1+q_k|\le\tau|r_k|$.
Since $D\ge0.0162079$, $|q_k|\le1/(0.0162079\,r_k^2)$, so $q_k<-1$ requires $|r_k|<1/\sqrt{0.0162079}<7.86$, and then $q_k\ge-1-7.86\tau$.
Summing $\text{acc}_k$ over the 50 stages, every $\tau$-feasible point has objective at least $-0.2185(1+7.86\tau)-50\tau$.
\end{remark}

% ---------------------------------------------------------------------------
\subsubsection{\inst{ex6_2_5} and \inst{ex6_2_7}}\label{app:small-ex62}

\paragraph{Models.}
In the notation of \cref{sec:other-dense}, $n_{ip}$ is the amount of component $i\in\{1,2,3\}$ in phase $p\in\{1,2,3\}$, with bounds $10^{-7}\le n_{ip}\le b_i$; the only rows are the three mass balances $\sum_p n_{ip}=b_i$, and $f(n)=\sum_{p}G_p(n_{\cdot p})$ with $n_{\cdot p}=(n_{1p},n_{2p},n_{3p})$.
In \inst{ex6_2_7} (ethylene glycol, lauryl alcohol, nitromethane) the three phases are liquid with the same UNIQUAC function, $b=(0.4,0.1,0.5)$ and $T:=\sum_ib_i=1$.
In \inst{ex6_2_5} (sec-butyl alcohol, di-sec-butyl ether, water) two liquid phases share a UNIQUAC function, $b=(40.30707,5.14979,54.54314)$, $T=100$, and the third phase is an ideal vapour with
\[
G_3(n)=\sum_i n_i\Bigl(\ln\frac{n_i}{\sum_k n_k}+0.156969560191053\Bigr).
\]
The OSIL objective is the expanded GAMS-Convert form of UNIQUAC (26 terms per phase in \inst{ex6_2_7}); the source models are Test Problems 5 and 7 of Chapter~6 of \citet{floudas1999-handbook-of-test-problems-in}.
After logarithms of quotients are split, each $G_p$ is a finite sum of \emph{linear atoms} $\alpha^\top n_{\cdot p}$
and \emph{log atoms} $(\ell^\top n_{\cdot p})\ln(m^\top n_{\cdot p})$ with $m\ge0$, $m\neq0$, so that $m^\top n_{\cdot p}>0$ for $n_{\cdot p}>0$.
No term mixes phases.
For each phase define the \emph{residual vector} $\rho_p:=\sum\ell$, the sum over the log atoms of $G_p$.
For the stored decimals, read as exact rationals,
\[
\rho_p=\bigl(0,\,0,\,\tfrac{1}{20000000000000}\bigr)\ \text{(each phase of \inst{ex6_2_7})},\qquad \rho_p=0\ \text{(each phase of \inst{ex6_2_5})}.
\]
The \inst{ex6_2_7} residual comes from three constants of the third component:
$8.73945638067505+1.868-10.607456380675=\sci{5}{-14}$.
Three separately written codes (a rule-based analysis of the expression trees, a symbolic identity check, and a pattern-matching decomposition that merges atoms with the same logarithm argument exactly) derived these atom decompositions and residuals and confirmed that the phase functions named equal above are identical expressions.

\begin{lemma}[Scaling identity]\label{lem:ex62-scaling}
For every $t>0$ and $n\in\R^3$ with $n>0$, $G_p(tn)=t\,G_p(n)+t\ln t\,\langle \rho_p,n\rangle$.
\end{lemma}

\begin{proof}
Linear atoms are homogeneous of degree one.
For a log atom, $m^\top n>0$ gives
\[
(\ell^\top tn)\ln(m^\top tn)=t\,(\ell^\top n)\ln(m^\top n)+t\ln t\,(\ell^\top n).
\]
Summing over the atoms gives the identity.
\end{proof}

\begin{lemma}[Ideal phase]\label{lem:ex62-ideal}
Let $G(y)=\sum_i y_i(\ln y_i+c)$ on $\{y>0:\sum_iy_i=1\}$, and let $\lambda\in\R^3$.
Then $G(y)-\lambda^\top y\ge-\ln\sum_ie^{\lambda_i-c}$.
\end{lemma}

\begin{proof}
Put $Z=\sum_ie^{\lambda_i-c}$ and $\pi_i=e^{\lambda_i-c}/Z$.
Then $G(y)-\lambda^\top y=\sum_iy_i\ln(y_i/\pi_i)-\ln Z\ge-\ln Z$ by Gibbs' inequality.
\end{proof}

On the simplex the vapour function of \inst{ex6_2_5} equals $\sum_i y_i(\ln y_i+0.156969560191053)$, because $\sum_k y_k=1$.

\begin{proposition}[Mass-balance Lagrangian]\label{prop:ex62-dual}
Let $\varepsilon=10^{-7}/T$ and $\Delta_\varepsilon=\{y\in\R^3:y_i\ge\varepsilon,\ \sum_iy_i=1\}$.
Let $\lambda\in\R^3$ be arbitrary, and let $m_p\le\inf_{y\in\Delta_\varepsilon}\bigl(G_p(y)-\lambda^\top y\bigr)$ for $p=1,2,3$.
If $\rho_p\ge0$ componentwise for every $p$, then every exactly feasible $n$ satisfies
\begin{equation}\label{eq:ex62-bound}
f(n)\ \ge\ \lambda^\top b+\sum_{p=1}^{3}\Bigl(\min(0,\,T m_p)-\frac{\max_i \rho_{p,i}}{e}\Bigr).
\end{equation}
\end{proposition}

\begin{proof}
The mass balances give $\sum_p\lambda^\top n_{\cdot p}=\lambda^\top b$, hence
$f(n)=\lambda^\top b+\sum_p\bigl(G_p(n_{\cdot p})-\lambda^\top n_{\cdot p}\bigr)$.
Fix $p$ and let $t=\sum_in_{ip}$ and $y=n_{\cdot p}/t$.
All amounts are at least $10^{-7}$, so $t>0$, and $n_{ip}\le\sum_qn_{iq}=b_i$ gives $t\le T$.
Hence $y_i\ge10^{-7}/t\ge\varepsilon$ and $y\in\Delta_\varepsilon$.
By \cref{lem:ex62-scaling},
$G_p(n_{\cdot p})-\lambda^\top n_{\cdot p}=t\bigl(G_p(y)-\lambda^\top y\bigr)+t\ln t\,\langle \rho_p,y\rangle$.
The first term is at least $t\,m_p\ge\min(0,Tm_p)$ because $0<t\le T$.
In the second term $\langle \rho_p,y\rangle\in[0,\max_i\rho_{p,i}]$ and $t\ln t\ge-1/e$, so the term is at least $-\max_i\rho_{p,i}/e$.
\end{proof}

The proof uses neither the upper bounds $n_{ip}\le b_i$ nor any property of $\lambda$.
If $\lambda$ are the chemical potentials of a stable equilibrium, $G_p(y)-\lambda^\top y$ is the classical tangent-plane distance of phase $p$, which is nonnegative on the simplex and vanishes at the equilibrium compositions; this is why the bounds $m_p$ below are close to $0$.

\paragraph{Computation.}
The multipliers are decimals with 20 significant digits, from a tangent-plane linear program on a grid of the simplex followed by a 60-digit Newton solve of the 12-equation KKT system;
how $\lambda$ was found does not affect validity.
\begin{align*}
\text{\inst{ex6_2_7}:}\quad \lambda_1&=-0.23993666802341555716, & \lambda_2&=-0.54068374800661316808,\\
\lambda_3&=-0.021609146907096643708, & \lambda^\top b&=-0.160847615463575861526;\\
\text{\inst{ex6_2_5}:}\quad \lambda_1&=-0.92114611232187116319, & \lambda_2&=-2.2777893037791928146,\\
\lambda_3&=-0.40139310690864392314, & \lambda^\top b&=-70.7520778334477055803539469 .
\end{align*}
The second code bounds $m_p$ by a two-dimensional interval branch and bound in the coordinates $(y_1,y_2)$, with $y_3=1-y_1-y_2$ enclosed on each box and clipped to $[y_{\min},1]$.
The domain is $[y_{\min},1-2y_{\min}]^2\cap\{y_1+y_2\le1-y_{\min}\}$ with $y_{\min}=\sci{9.9}{-8}$ (\inst{ex6_2_7}) and $\sci{9.9}{-10}$ (\inst{ex6_2_5}), slightly below $\varepsilon$, so the domain contains $\Delta_\varepsilon$.
A box is discarded without bounding only if the exact test $\underline{y}_1+\underline{y}_2>1-y_{\min}$ shows that it misses the domain.
Write $\psi(y_1,y_2)$ for $G_p(y)-\lambda^\top y$ with $y_3=1-y_1-y_2$.
On each box the lower bound is the best of the natural interval extension, the mean-value form at a feasible centre $c$, and the following second-order form.
Let $[\underline H,\overline H]$ be an interval enclosure of the Hessian of $\psi$ over the box, and put $a=\underline H_{11}$, $d=\underline H_{22}$,
\[
Q^-=\begin{pmatrix}a&\underline H_{12}\\ \underline H_{12}&d\end{pmatrix},\qquad Q^+=\begin{pmatrix}a&\overline H_{12}\\ \overline H_{12}&d\end{pmatrix}.
\]
For every displacement $z$ and every symmetric $H$ in the enclosure, $z^\top Hz$ is at least the smaller of $z^\top Q^-z$ and $z^\top Q^+z$: the diagonal terms are bounded below independently, since $z_i^2\ge0$, and the off-diagonal term $2H_{12}z_1z_2$ is linear in $H_{12}$, so it is minimal at an endpoint.
By Taylor's theorem, with the Hessian at a point of the segment from $c$ to $c+z$, every $c+z$ in the box satisfies $\psi(c+z)\ge\psi(c)+\nabla\psi(c)^\top z+\tfrac12\min(z^\top Q^-z,\,z^\top Q^+z)$.
When $Q^-$ and $Q^+$ are positive definite, the smaller of the two unconstrained quadratic minima, $\psi(c)-\tfrac12\max_\pm\nabla\psi(c)^\top(Q^\pm)^{-1}\nabla\psi(c)$, is therefore a valid lower bound on the box; otherwise the quadratic term with the lower diagonal ends is enclosed over the box in interval arithmetic.
The code evaluates these closed-form minima in outward-rounded mpmath interval arithmetic, with interval enclosures of $\psi(c)$ and $\nabla\psi(c)$.
A box is fathomed when its lower bound is at least $-\tau$, so $m_p=-\tau$ is valid.
The arithmetic is mpmath interval arithmetic with 53-bit boxes and 40-digit centre values, applied to symbolic derivatives of the OSIL trees.
For \inst{ex6_2_7}, $\tau=\sci{6}{-15}$ (42{,}111 boxes, one run for the phase function shared by the three phases);
for \inst{ex6_2_5}, $\tau=10^{-17}$ (131{,}111 boxes for the liquid function), and for the vapour phase \cref{lem:ex62-ideal} gives $m_3\ge\sci{6.978}{-22}>0$, so its term in \eqref{eq:ex62-bound} is $0$.
The assembly of \eqref{eq:ex62-bound} is exact rational arithmetic, with a rational lower bound on $e$:
\begin{align*}
\text{\inst{ex6_2_7}:}\quad & \lambda^\top b-3\cdot\sci{6}{-15}-3\cdot\sci{5}{-14}/e\ >\ -0.16084761546364905,\\
\text{\inst{ex6_2_5}:}\quad & \lambda^\top b-2\cdot100\cdot10^{-17}=-70.7520778334477075803539469\\
&\hphantom{\lambda^\top b-2\cdot100\cdot10^{-17}}>-70.75207783344770759.
\end{align*}
The right sides are the displayed values of $\Lcert$.
In the \inst{ex6_2_7} bound the residual term $3\cdot\sci{5}{-14}/e$ contributes $\sci{5.52}{-14}$, more than the final gap.

\paragraph{Why $\lambda^\top b$ alone is not a bound.}
For neither instance is $\lambda^\top b$ a valid dual bound: the exactly feasible points below lie \sci{2.5}{-14} (\inst{ex6_2_7}) and \sci{1.3}{-20} (\inst{ex6_2_5}) below it.
For \inst{ex6_2_7} the reason is the residual $\rho_p$:
with this $\lambda$ the tangent-plane distance has a local minimum of about $\sci{-4.19}{-15}$, so no $\tau$ below about \sci{4.19}{-15} can be certified, and a run with $\tau=10^{-17}$ failed as expected.
For \inst{ex6_2_5} the reason is the rounding of $\lambda$: the local minima of the liquid distance are about $\sci{-1.24}{-22}$ and $\sci{-1.38}{-21}$.

\paragraph{Exactly feasible points.}
Each point is rational: the amounts of two phases are 20-digit decimals of the KKT amounts, and the third phase is $b$ minus their sum, computed exactly.
The three rows then hold exactly, and all bounds were checked exactly.
The objectives were enclosed twice:
by mpmath interval arithmetic, and by a separate evaluation, without any library function, in rational interval arithmetic in which each logarithm is enclosed by its $\operatorname{atanh}$ series with a remainder bound and outward rounding to $2^{-400}$.
The second evaluation gives
\begin{align*}
f(n^\star_{\text{\inst{ex6_2_7}}})&\in[-0.1608476154636008615244708,\ -0.1608476154636008615244707],\\
f(n^\star_{\text{\inst{ex6_2_5}}})&\in[-70.7520778334477055803671221,\ -70.7520778334477055803671220].
\end{align*}

\paragraph{Proof of \cref{thm:ex62-bounds}.}
The residual vectors are nonnegative, so \cref{prop:ex62-dual} applies with the multipliers and tangent-plane minima above and gives the two values of $\Lcert$.
The rational points above are exactly feasible, and the upper ends of their objective enclosures give
$\Uprim=-0.16084761546360086$ and $\Uprim=-70.75207783344770558$ (rounded up).
From the certificate ends, $\gapabs\le\sci{4.82}{-14}$ for \inst{ex6_2_7} and $\gapabs\le\sci{2.00}{-15}$ for \inst{ex6_2_5}, that is $\gaprel\le\sci{3.00}{-13}$ and $\gaprel\le\sci{2.83}{-17}$.
For \inst{ex6_2_5} the printed displays differ by $\sci{2.01}{-15}$, slightly more than the gap cell, which uses the certificate ends (\cref{tab:closures}).
\hfill$\square$

\begin{certbox}[Certificate for \cref{thm:ex62-bounds}]
\certitem{Inputs}{OSIL files (SHA-256 \texttt{eed48057}, \texttt{c7f7cd37}); 20-digit multipliers; rational primal points.}
\certitem{Computation}{boxes \arith{I} ($\log$); assembly \arith{E}; primal rows \arith{E}, objectives \arith{I} and library-free rational intervals.}
\certitem{Data reading}{exact, then outward in the interval steps.}
\certitem{Implementations}{displayed bounds from the second code. The first code (binary64 intervals widened by one ulp per operation, its own rigorous logarithm, other coordinates and splitting rule, $\tau=10^{-11}$) certifies the weaker $-0.16084761549352554$ (\inst{ex6_2_7}) and $-70.752077836333563$ (\inst{ex6_2_5}). A third separately written code (mpmath intervals and exact rationals, own OSIL reader, closed-form derivatives) proves convexity windows around the local minimizers of the tangent-plane distance, excludes the rest of the simplex by interval branch and bound and certifies slightly stronger bounds; a separate agent session read it and reran \inst{ex6_2_7}; two of its consistency assertions are incomplete (the sign of $\rho_p$ and the form of the vapour function), and the other two codes establish both properties. Two evaluations enclose the primal objectives.}
\certitem{Evidence level}{\evid{rerun} (the search trees are not stored).}
\certitem{Replay}{Tier~1; 258\,s and 77\,s (198\,s and 63\,s for the third code).}
\end{certbox}

% ---------------------------------------------------------------------------
\subsubsection{\inst{etamac}}\label{app:small-etamac}

\paragraph{Model.}
The instance is the GAMS Model Library model \texttt{etamac}, an energy--economy model after the ETA-MACRO model of \citet{manne1977-eta-macro-a-model-of}.
It has 97 variables and 70 rows: 60 linear equalities, 9 nonlinear equalities and 1 linear inequality.
There are nine periods $t=1,\dots,9$.
We write $K,K\!N,Y,Y\!N,C,I,E\!C$ for capital, new capital, output, new output, consumption, investment and energy cost,
and $L,E,L\!N,E\!N$ for the two priced inputs of the second nest and their new vintages (our labels).
Put $g=4.91287681$ and $\delta=0.8153726976$.
The objective is $\min f=-\sum_t\beta_t\ln C_t$ with decimal weights $\beta_t>0$.
The linear rows are
\begin{gather*}
K\!N_t=g\,I_{t-1},\qquad K_t=\delta K_{t-1}+K\!N_t,\qquad Y_t=\delta Y_{t-1}+Y\!N_t\qquad(t\ge2),\\
L_1=L\!N_1+2.038431744,\qquad L_t=\delta L_{t-1}+L\!N_t\qquad(t\ge2),\\
E_1=E\!N_1+40.76863488,\qquad E_t=\delta E_{t-1}+E\!N_t\qquad(t\ge2),\\
1000\,E\!C_t=c^L_tL_t+c^E_tE_t\ \ (c^L_t,c^E_t>0),\qquad Y_t=C_t+I_t+E\!C_t,\qquad 0.07\,K_9\le I_9 .
\end{gather*}
The production rows are $Y\!N_t=\Phi_t(K\!N_t,L\!N_t,E\!N_t)$ for $t\ge2$ and $Y_1=3.4653339648+\Phi_1(L\!N_1,E\!N_1)$, where
\[
\Phi_t=\bigl(a_t\,K\!N_t^{-p_1}+b\,L\!N_t^{-p_2}E\!N_t^{-p_3}\bigr)^{-q},\qquad
\Phi_1=\bigl(c_0+b\,L\!N_1^{-p_2}E\!N_1^{-p_3}\bigr)^{-q},
\]
with $a_t\in[0.330,0.821]$ and
\begin{gather*}
p_1=0.342222222222222,\quad p_2=0.427777777777778,\quad p_3=0.794444444444445,\\
q=0.818181818181818,\quad b=0.306708090151268,\quad c_0=0.612508399277048.
\end{gather*}
$K_1=12.32657617084$ is fixed, $E\!C_t$ is free, and every other variable has a positive lower bound and no upper bound.

The objective is convex and all other rows are linear, so if every $\Phi_t$ were concave, relaxing the production equalities to ``$\le$'' would give a convex program (\cref{sec:other-convex}).
The stored exponents, however, are 15-digit roundings of $0.28\cdot11/9$, $0.35\cdot11/9$, $0.65\cdot11/9$ and $9/11$.
In exact rationals $p_1q=0.2799999999999997559\ldots$, and the degree of the Cobb--Douglas aggregate $L\!N^{p_2q}E\!N^{p_3q}$ is
\[
s:=(p_2+p_3)\,q=1+\frac{207070707070707}{5\cdot10^{29}}\approx1+\sci{4.14}{-16}.
\]
So $\Phi_t$ is not concave on the positive orthant, and the certificate must repair this.

\begin{lemma}[A box from the rows]\label{lem:etamac-box}
There are rationals $lo_j\le hi_j$ such that every exactly feasible point lies in $\mathcal B=\prod_j[lo_j,hi_j]$.
The $lo_j$ are the OSIL lower bounds, and $lo=(c^L_t\,lo_{L_t}+c^E_t\,lo_{E_t})/1000$ for $E\!C_t$.
\end{lemma}

\begin{proof}
The upper ends follow by a forward recursion over $t$, each step using one row and bounds already established:
$Y_1\le3.4653339648+c_0^{-q}$, since the bracket of $\Phi_1$ is at least $c_0$ and $x\mapsto x^{-q}$ is decreasing;
$I_t\le Y_t-lo_{C_t}-lo_{E\!C_t}$ from $Y_t=C_t+I_t+E\!C_t$;
$K\!N_{t+1}=g\,I_t$;
$Y\!N_{t+1}\le a_{t+1}^{-q}K\!N_{t+1}^{p_1q}$, by dropping the positive $L\!N$--$E\!N$ term inside $\Phi_{t+1}$;
$Y_{t+1}=\delta Y_t+Y\!N_{t+1}$;
and the bounds of $C$, $E\!C$, $L$, $E$, $L\!N$, $E\!N$ and $K$ from the linear rows with positive coefficients.
The recursion is evaluated in outward-rounded interval arithmetic.
\end{proof}

The largest upper end is $2062.83$ ($E_9$); for example $Y_1\le4.959$ and $Y_9\le26.73$.

\begin{lemma}[Concave majorant]\label{lem:etamac-majorant}
Let $w=L\!N^{p_2/(p_2+p_3)}E\!N^{p_3/(p_2+p_3)}$, $W\ge\max_{\mathcal B}w$ and $\kappa\ge\max(1,W^{s-1})$.
Define
\[
\tilde\Phi_t=\bigl(a_t\,K\!N^{-p_1}+b\,\kappa^{-1/q}L\!N^{-p_2/s}E\!N^{-p_3/s}\bigr)^{-q},
\]
and $\tilde\Phi_1$ likewise with $c_0$ in place of $a_tKN^{-p_1}$.
Then (a) $\tilde\Phi_t$ is concave and nondecreasing on the positive orthant, and (b) $\Phi_t\le\tilde\Phi_t$ on $\mathcal B$.
\end{lemma}

\begin{proof}
For $\alpha,\beta>0$, $M(u,v)=(\alpha u^{-1/q}+\beta v^{-1/q})^{-q}$ is a positive multiple of a weighted power mean with exponent $-1/q<1$,
hence concave and nondecreasing in $(u,v)$ on the positive orthant.
Put $u=K\!N^{p_1q}$.
Then $u^{-1/q}=K\!N^{-p_1}$, $(w^s)^{-1/q}=L\!N^{-p_2}E\!N^{-p_3}$ and $(\kappa w)^{-1/q}=\kappa^{-1/q}L\!N^{-p_2/s}E\!N^{-p_3/s}$,
so $\Phi_t=M(u,w^s)$ and $\tilde\Phi_t=M(u,\kappa w)$ with $\alpha=a_t$, $\beta=b$.
The function $u$ is concave because $0<p_1q<1$, and $w$ is concave because it is a Cobb--Douglas function with exponents summing to one.
A concave function that is nondecreasing in each argument, composed with concave functions, is concave; this gives (a).
For (b), $s>1$ and $0<w\le W$ give $w^s=w\cdot w^{s-1}\le\kappa w$, and $M$ is nondecreasing in $v$.
For $\tilde\Phi_1$ the first argument is replaced by a constant.
\end{proof}

Let $R$ be the problem obtained from \inst{etamac} by adding $x\in\mathcal B$ and replacing $Y\!N_t=\Phi_t$ by $Y\!N_t\le\tilde\Phi_t$ ($t\ge2$) and the row for $Y_1$ by $Y_1\le3.4653339648+\tilde\Phi_1$.
By \cref{lem:etamac-box,lem:etamac-majorant}(b), every exactly feasible point of \inst{etamac} is feasible for $R$,
and by \cref{lem:etamac-majorant}(a), $R$ is a convex program.

\begin{proposition}[Tangent-plane bound]\label{prop:etamac-tangent}
Let $h(x)=0$ collect the 60 linear equalities, $\tilde g(x)\le0$ the 9 relaxed production rows and $g_{70}(x)\le0$ the terminal row.
For any $\nu\in\R^{60}$, any $\mu\in\R^{10}$ with $\mu\ge0$ and any $\hat x\in\mathcal B$, put $l=f+\nu^\top h+\mu^\top(\tilde g,g_{70})$.
Then every exactly feasible $x$ satisfies
\[
f(x)\ \ge\ l(\hat x)+\sum_j\ \min_{\xi\in[lo_j,hi_j]}\ \partial_jl(\hat x)\,(\xi-\hat x_j).
\]
\end{proposition}

\begin{proof}
A feasible $x$ lies in $\mathcal B$ and in $R$, so $h(x)=0$, $\tilde g(x)\le0$ and $g_{70}(x)\le0$; with $\mu\ge0$ this gives $f(x)\ge l(x)$.
On a neighbourhood of $\mathcal B$ (all variables except $E\!C$ are positive there, and $E\!C$ enters linearly) $l$ is convex and differentiable:
$f$ is convex, $h$ and $g_{70}$ are affine, $\tilde g$ is convex by \cref{lem:etamac-majorant}(a), and $\mu\ge0$.
Hence $l(x)\ge l(\hat x)+\nabla l(\hat x)^\top(x-\hat x)$, and minimizing the right side coordinatewise over $\mathcal B$ gives the claim.
\end{proof}

\paragraph{Computation and proof of \cref{thm:etamac-bound}.}
The second code matched all 70 rows, the 97 bounds and the objective to templates built from an explicit map of variable roles.
It computed $\mathcal B$ by \cref{lem:etamac-box}, $W=1015.6000321087411$ (the maximum over $t$ of $w$ at the upper ends of $L\!N_t$ and $E\!N_t$),
$W^{s-1}\le1+\sci{2.8672}{-15}$ and $\kappa=1.000000000000004$.
The point $\hat x$ and the multipliers $(\nu,\mu)$ come from a local solve of $R$ followed by a 50-digit Newton solve of the KKT system of $R$ (residual \sci{2.7}{-48}), rounded to 30 digits.
The multipliers $\mu$ are positive (the smallest is $0.277935$, checked on the decimal strings), and every component of $\hat x$ that is not fixed lies at least $0.047$ above its lower bound.
The value $l(\hat x)$ and the gradient $\nabla l(\hat x)$ are evaluated in 40-digit mpmath interval arithmetic from symbolic expressions with exact rational data.
The largest gradient component of a variable that is not fixed is \sci{4.9}{-31};
the component for the fixed $K_1$ is $0.00496$ and multiplies a range of width zero.
\Cref{prop:etamac-tangent} then gives $f\ge-15.2946756433680921685$ (rounded down), displayed as $\Lcert=-15.29467564336809217$.

The exactly feasible point is defined from the first code's saved 17-digit decisions $I_1,\dots,I_8$, $L\!N_t$ and $E\!N_t$ by forward definitions (\cref{prop:pts-triangular}):
$I_9:=0.07\,K_9$, so the terminal row is active, and every other variable is defined by the row that determines it from earlier ones.
The rows therefore hold as real identities.
Interval evaluation at 40 or more digits, in one implementation, shows that every nonfixed variable satisfies its finite bounds with margin at least $0.047$ and that $K_1$ equals its fixed value exactly,
and encloses the objective in an interval of width below $10^{-44}$ with upper end $-15.29467564336808959198\ldots$, displayed as $\Uprim=-15.29467564336808959$.
The certified gap is $\gapabs\le\sci{2.58}{-15}$ (exact value $\sci{2.5765}{-15}$), and $\gaprel\le\sci{1.69}{-16}$.
\hfill$\square$

The certified bound differs from $l(\hat x)$ only by the gradient term; we did not separate the contributions of $\kappa$, of the degree excess $s-1$ and of the suboptimality of the point to the remaining gap.

\begin{certbox}[Certificate for \cref{thm:etamac-bound}]
\certitem{Inputs}{OSIL file (SHA-256 \texttt{6716cc88}); 30-digit KKT point and multipliers; 17-digit primal decisions.}
\certitem{Computation}{$\mathcal B$, $\kappa$, $l(\hat x)$ and $\nabla l(\hat x)$ at 40 digits, \arith{I} ($\exp$, $\log$); primal point \arith{I}.}
\certitem{Data reading}{exponents exact; other data exact or outward; rows, bounds and objective matched to templates.}
\certitem{Implementations}{the second code certifies the displayed bound (a later separate agent session read it line by line and found no error); the first code, linearized at the KKT point of the unrelaxed model with the cruder $W=2022.06$, certifies the weaker $-15.294675643368096$; primal enclosure by one code.}
\certitem{Evidence level}{\evid{stored}.}
\certitem{Replay}{Tier~1; 22\,s (27\,s for the first code).}
\end{certbox}

% ---------------------------------------------------------------------------
\subsubsection{\inst{pricing050}}\label{app:small-pricing}

\paragraph{Model.}
The instance is a continuous version of the pricing model of Davarnia and van Hoeve with 50 products.
It maximizes, so in this subsection the dual bound $\Lcert$ is an upper bound and the directions of \cref{sec:semantics-display} are reversed.
\Cref{sec:other-dense} displays the problem; its rows are e2--e6, and its data are
integers $c_j\in\{0,\dots,20\}$ (46 of them positive), one-decimal coefficients $a_{ij}<0$, exponents $p_{ij}\in\{1,2,3\}$ (78, 94 and 77 terms) and
$r=(-500,-651,-615,-788,-984)$.
The coefficients $g_{ij}$ are $-0.1$, $-1.0000000000000002\mathrm{e}{-2}$ and $-1.0000000000000002\mathrm{e}{-3}$ for $p_{ij}=1,2,3$:
the last two are the shortest round-trip decimals of the binary64 products $\fl(0.1\cdot0.1)$ and $\fl(0.1^3)$, and the model uses these decimals exactly.
Integer powers are exact, also at $x_j=0$ (\cref{sec:semantics-model}).

Write $\mathrm{row}_i(x)=\sum_ja_{ij}x_j\exp(g_{ij}x_j^{p_{ij}})$ and
$F_j(\xi)=c_j\xi+\sum_i\mu_ia_{ij}\,\xi\,e^{g_{ij}\xi^{p_{ij}}}$.

\begin{proof}[Proof of \eqref{eq:pricing-lagrangian}]
Let $\mu\in\R^5$ with $\mu\ge0$, and let $x$ be exactly feasible.
Then $\mu_i(\mathrm{row}_i(x)-r_i)\le0$ for every $i$, so
$c^\top x\ge c^\top x+\sum_i\mu_i(\mathrm{row}_i(x)-r_i)=\sum_jF_j(x_j)-\mu^\top r\ge\sum_j\min_{\xi\in[0,10]}F_j(\xi)-\mu^\top r$.
Multiplying by $-1$ gives \eqref{eq:pricing-lagrangian}.
\end{proof}

\paragraph{Computation and proof of \cref{thm:pricing-bound}.}
The multipliers are
\[\mu_{e5}=3.0489011208166370021,\qquad\mu_{e6}=2.1677509642745686136,\qquad\mu_{e2}=\mu_{e3}=\mu_{e4}=0;\]
exactly, $\mu^\top r=-4535.6010320496854734372$.
The second code partitions each $[0,10]$ into pieces on which interval arithmetic proves $F_j'>0$, $F_j'<0$ or $F_j''<0$, or proves $F_j''>0$ on a piece of width below $10^{-7}$ and then uses an exact quadratic bound (2{,}406 pieces in total), and obtains $\sum_j\min F_j\ge-2721.771953597712415687255$.
A separate check computes a 50-digit stationary point of each $F_j$, proves $F_j''>0$ on a window around it (so $F_j\ge F_j(\hat\xi)-F_j'(\hat\xi)^2/(2\underline{F_j''})$ there), and excludes the rest of $[0,10]$ by interval branch and bound (247 boxes), obtaining $\sum_j\min F_j\ge-2721.77195359771241568725405377$.
Then \eqref{eq:pricing-lagrangian} gives $-c^\top x\le-1813.82907845197305774994594623$, displayed rounded up as $\Lcert=-1813.8290784519730577$.

The exactly feasible point is the first code's saved KKT point with 17 significant digits, in which $x_{11}$ was raised by \sci{1.27}{-15}.
All bounds hold exactly, and interval evaluation shows that rows e5 and e6 hold with slack at least \sci{3.4958}{-15} and \sci{3.9173}{-15}, and rows e2--e4 with slack above $76$.
Two separately written implementations, each with its own OSIL reader, check this point and find the same slacks and the same exact objective: the earlier separate check of the certificate box below and the separate check above, whose session had not read the earlier check.
Because the $c_j$ are integers and the point is decimal, its objective is exactly $\Uprim=-1813.8290784519730769$.
So $\gapabs\le\sci{1.92}{-14}$ (exact value $\sci{1.9150}{-14}$) and $\gaprel\le\sci{1.06}{-17}$.
\hfill$\square$

\begin{certbox}[Certificate for \cref{thm:pricing-bound}]
\certitem{Inputs}{OSIL file (SHA-256 \texttt{56d3e00f}); 20-digit multipliers; 17-digit primal point.}
\certitem{Computation}{50 univariate minima \arith{I} ($\exp$); assembly \arith{E}; row slacks \arith{I}, objective \arith{E}.}
\certitem{Data reading}{exact, then outward in the interval steps.}
\certitem{Implementations}{the second code and a separate check reproduce the displayed bound; an earlier separate check with its own reader differs from the separate check by about \sci{2}{-27}; the first code (interval branch and bound, 16{,}960 boxes) certifies a weaker one; the primal point is checked by the earlier separate check and by the separate check, two separately written codes.}
\certitem{Evidence level}{\evid{stored}.}
\certitem{Replay}{Tier~1; 28\,s (3\,s for the separate check).}
\end{certbox}

% ---------------------------------------------------------------------------
\subsubsection{\inst{pindyck}}\label{app:small-pindyck}

\paragraph{Model.}
The instance is the GAMS Model Library model \texttt{pindyck} (OPEC pricing and extraction, after \citealp{pindyck1978-gains-to-producers-from-the}).
It has 116 variables, $7\cdot16$ per-period variables plus four fixed initial states, and 96 equality rows.
For $t=1,\dots,16$, with prices $p_t$,
\begin{gather*}
td_t=0.87\,td_{t-1}-0.13\,p_t+c_t,\qquad s_t=0.75\,s_{t-1}+1.02^{-k\,cs_t}\,(1.1+0.1\,p_t),\qquad cs_t=cs_{t-1}+s_t,\\
d_t=td_t-s_t,\qquad R_t=R_{t-1}-d_t,\qquad \pi_t=(p_t-250/R_t)\,d_t,
\end{gather*}
with $td_0=18$, $s_0=6.5$, $cs_0=0$, $R_0=500$, exogenous decimals $c_t$ increasing from $3.3$ to $3.87553375330505$, and $k=0.142857142857143$.
The objective is $\min-\sum_t\delta_t\pi_t$, where $\delta_t$ are the 15-digit decimals of $1.05^{-(t-1)}$.
All variables except $\pi_t$ are nonnegative.
Write $K=k\ln1.02>0$, so that $1.02^{-k\,cs_t}=e^{-K\,cs_t}$.

\begin{lemma}[Reduction to prices]\label{lem:pindyck-reduction}
For every $p\in\R^{16}$ with $p\ge0$ the recursions have a unique solution $td(p),s(p),cs(p),d(p),R(p)$, and where every $R_t(p)\ne0$ the rows also determine $\pi(p)$.
Let $P=\{p\ge0:d_t(p)\ge0\ \forall t\}$.
Every exactly feasible point has prices in $P$ and objective $-J(p)$, where $J(p)=\sum_t\delta_t\,d_t(p)\,(p_t-250/R_t(p))$.
\end{lemma}

\begin{proof}
Given $p$ and the states of period $t-1$, the supply row reads $\psi(s):=s-a-b\,e^{-Ks}=0$ with $a=0.75\,s_{t-1}$ and $b=(1.1+0.1\,p_t)\,e^{-K\,cs_{t-1}}>0$.
The function $\psi$ is continuous and strictly increasing with limits $\mp\infty$, so it has exactly one root.
The other rows define the remaining states explicitly.
A feasible point satisfies $d_t\ge0$ and $R_t\neq0$, and its objective is $-J(p)$.
The set $P$ drops the sign conditions on $td$, $s$, $cs$ and $R$, so it contains the prices of every feasible point; this is the safe direction for a lower bound.
\end{proof}

So it suffices to bound $J$ from above on $P$.
Sampled Hessians of $J$ are negative definite, but $P$ is not known to be convex, and interval Hessians computed entry by entry over the price box lose the correlations between entries.
The certificate therefore proves strong concavity on a polytope $G\supseteq P$ through an explicit map from per-period quantities to the Hessian.

\begin{lemma}[A polytope containing $P$]\label{lem:pindyck-polytope}
Let $\alpha_t=td_t(0)$, let $CSH_t\ge\sup_{p\in P}cs_t(p)$, and let $e_{t}^{\mathrm{lo}}\le\exp(-K\,CSH_t)$ be rationals.
For $j\le t$ put $W_{tj}=0.13\cdot0.87^{t-j}+0.1\cdot0.75^{t-j}e_j^{\mathrm{lo}}$ and
$\gamma_t=\alpha_t-6.5\cdot0.75^t-1.1\sum_{j\le t}0.75^{t-j}e_j^{\mathrm{lo}}$.
Let $\underline W\le W$ and $\bar\gamma\ge\gamma$ entrywise, and $G=\{p\ge0:\underline Wp\le\bar\gamma\}$.
Then $P\subseteq G$.
\end{lemma}

\begin{proof}
For $p\in P$, $e^{-K\,cs_t}\ge e_t^{\mathrm{lo}}$, so by induction
$s_t\ge\ell_t(p):=6.5\cdot0.75^t+\sum_{j\le t}0.75^{t-j}(1.1+0.1\,p_j)\,e_j^{\mathrm{lo}}$.
The condition $d_t\ge0$ gives $td_t(p)\ge s_t\ge\ell_t(p)$.
Since $td_t(p)=\alpha_t-0.13\sum_{j\le t}0.87^{t-j}p_j$, this is $(Wp)_t\le\gamma_t$.
Rounding is safe because $p\ge0$.
\end{proof}

The bounds $CSH_t$ come from six rounds of linear programs in $(p,s)$ whose feasible sets contain $(p,s(p))$ for every $p\in P$:
the supply row is relaxed by enclosing $e^{-K\,cs_t}$ between its values at the current bounds of $cs_t$,
and the programs use $s_t\le td_t(p)$, $p_t\le\alpha_t/0.13$ and $s_t\le\alpha_t$, which hold on $P$.
The first round starts from $cs_t\in[0,\sum_{j\le t}\alpha_j]$.
Each optimal value is certified by weak duality in exact rational arithmetic with finite variable bounds; the LP solver only proposes multipliers \citep{neumaier2004-safe-bounds-in-linear-and}.
The result is $CSH_{16}=169.273$.
On $G$ the prices satisfy $p_t\le\bar p_t:=\bar\gamma_t/\underline W_{tt}$, with $\bar p_t$ between $57.41$ and $126.15$.

\begin{lemma}[State ranges over $G$]\label{lem:pindyck-ranges}
There is a box $\Theta\subset\R^{112}$ that contains $\theta(p)$ for every $p\in G$, where $\theta$ is defined in \cref{lem:pindyck-hessian}.
On $G$ every $R_t$ is at least $192.5$.
\end{lemma}

\begin{proof}
Linear programs over $G$ in the variables $(p,s,\omega,z)$, where $\omega_t$ stands for $e^{-K\,cs_t}$ and $z_t$ for $p_t\omega_t$, contain $(p,s(p),\omega(p),z(p))$ for every $p\in G$.
They use the exact supply equation $s_t=0.75\,s_{t-1}+1.1\,\omega_t+0.1\,z_t$,
five tangent lines below and one secant above the convex function $c\mapsto e^{-Kc}$ on the current range of $cs_t$ (each line proved valid on that range in interval arithmetic),
the McCormick inequalities for $z_t=p_t\omega_t$,
and finite bounds on all variables, for example $s_t\le0.75\,s_{t-1}+1.1+0.1\,\bar p_t$ because $\omega_t\le1$.
Seven rounds, each certified by exact weak duality, give $cs_{16}\in[71.1,158.8]$; further certified programs give the ranges of $s_t$, $d_t$ and $R_t$,
for example $R_{16}\in[192.5,500.9]$ and $d_{16}\in[-3.16,23.48]$ ($d_t$ may be negative on $G$).
Interval arithmetic on these ranges gives $\Theta$.
\end{proof}

\begin{lemma}[Hessian map]\label{lem:pindyck-hessian}
For $p$ in a neighbourhood of $G$ put
$E_t=e^{-K\,cs_{t-1}}$, $\beta_t=(1.1+0.1\,p_t)E_t$, $\phi_t=e^{-K s_t}$, $u_t=p_t-250/R_t$,
$\theta_t=(\beta_t,E_t,\phi_t,d_t,u_t,R_t^{-2},R_t^{-3})$,
$\iota_t=1/(1+K\beta_t\phi_t)$, $\chi_t=\phi_t\iota_t$ and $\Gamma_t=\nabla cs_{t-1}$, with $\mathbf e_t$ the $t$-th unit vector.
Then $\nabla^2J(p)=\Psi(\theta(p))$, where $\Psi$ is given by the recursion
\begin{align*}
\nabla b_t&=0.1\,E_t\mathbf e_t-K\beta_t\Gamma_t,\\
\nabla^2b_t&=\beta_t\bigl(K^2\Gamma_t\Gamma_t^\top-K\nabla^2cs_{t-1}\bigr)-0.1\,K E_t\bigl(\mathbf e_t\Gamma_t^\top+\Gamma_t\mathbf e_t^\top\bigr),\\
\nabla s_t&=0.75\,\iota_t\nabla s_{t-1}+\chi_t\nabla b_t,\\
\nabla^2s_t&=0.75\,\iota_t\nabla^2s_{t-1}+\chi_t\nabla^2b_t-K\chi_t\bigl(\nabla b_t\nabla s_t^\top+\nabla s_t\nabla b_t^\top\bigr)+K^2\beta_t\chi_t\nabla s_t\nabla s_t^\top,\\
\nabla cs_t&=\nabla cs_{t-1}+\nabla s_t,\quad \nabla^2cs_t=\nabla^2cs_{t-1}+\nabla^2s_t,\quad
\nabla td_t=-0.13\textstyle\sum_{j\le t}0.87^{t-j}\mathbf e_j,\\
\nabla d_t&=\nabla td_t-\nabla s_t,\quad \nabla^2d_t=-\nabla^2s_t,\quad \nabla R_t=\nabla R_{t-1}-\nabla d_t,\quad \nabla^2R_t=\nabla^2R_{t-1}-\nabla^2d_t,\\
\nabla u_t&=\mathbf e_t+250R_t^{-2}\nabla R_t,\quad \nabla^2u_t=250\bigl(R_t^{-2}\nabla^2R_t-2R_t^{-3}\nabla R_t\nabla R_t^\top\bigr),\\
\nabla^2J&=\textstyle\sum_t\delta_t\bigl(u_t\nabla^2d_t+d_t\nabla^2u_t+\nabla d_t\nabla u_t^\top+\nabla u_t\nabla d_t^\top\bigr),
\end{align*}
started from $\nabla s_0=\nabla cs_0=\nabla R_0=0$ and zero Hessians.
\end{lemma}

\begin{proof}
The supply row is $s_t=a_t+b_t\phi_t$ with $a_t=0.75\,s_{t-1}$, $b_t=\beta_t$ and $\phi_t=e^{-Ks_t}$.
Differentiating once gives $(1+Kb_t\phi_t)\nabla s_t=\nabla a_t+\phi_t\nabla b_t$, which is the formula for $\nabla s_t$.
Differentiating $(1+Kb_t\phi_t)\,\partial_is_t=\partial_ia_t+\phi_t\,\partial_ib_t$ with respect to $p_j$ and using $\partial_j\phi_t=-K\phi_t\,\partial_js_t$ gives
\[
(1+Kb_t\phi_t)\,\partial_{ij}s_t=\partial_{ij}a_t+\phi_t\,\partial_{ij}b_t-K\phi_t\bigl(\partial_ib_t\,\partial_js_t+\partial_jb_t\,\partial_is_t\bigr)+K^2b_t\phi_t\,\partial_is_t\,\partial_js_t,
\]
which is the formula for $\nabla^2s_t$.
With $\nabla E_t=-KE_t\Gamma_t$, the product rule gives $\nabla b_t$ and $\nabla^2b_t$; the remaining lines are the product rule and linearity.
The functions are $C^2$ near $G$ by the implicit function theorem, since $1+Kb_t\phi_t>0$, and because $R_t\ge192.5>0$ on $G$ (\cref{lem:pindyck-ranges}).
\end{proof}

The recursion was derived twice without computer algebra, in the sessions that wrote the first and the second code; the implicit step was also checked symbolically,
and $\Psi$ agrees with 50-digit finite-difference Hessians at 12 points of $G$ to \sci{1.7}{-16}; the recursion as a whole has no computer-algebra proof.

\begin{lemma}[Matrix test]\label{lem:pindyck-matrix}
Let $\Theta'$ be a box, and suppose that for every $\theta\in\Theta'$,
$\Psi(\theta)=C+\sum_{k=1}^m\varepsilon_kA_k+E$ with $\varepsilon\in[-1,1]^m$, symmetric matrices $C$, $A_k$, $E$ (where $\varepsilon$ and $E$ may depend on $\theta$), and $|E|\le\bar E$ entrywise for a fixed matrix $\bar E$.
If $X_k\succeq\pm A_k$ for all $k$, $\bar\rho\ge\rho(\bar E)$ and $-C-\sum_kX_k-(\bar\rho+\mu)I\succ0$, then $\Psi(\theta)\preceq-\mu I$ on $\Theta'$.
\end{lemma}

\begin{proof}
For $|\varepsilon_k|\le1$, $X_k-\varepsilon_kA_k=\tfrac12(1+\varepsilon_k)(X_k-A_k)+\tfrac12(1-\varepsilon_k)(X_k+A_k)\succeq0$.
For a unit vector $v$, $v^\top Ev\le|v|^\top|E|\,|v|\le\rho(|E|)\le\rho(\bar E)\le\bar\rho$, by the Perron--Frobenius monotonicity of the spectral radius on nonnegative matrices.
Hence $\Psi\preceq C+\sum_kX_k+\bar\rho I\prec-\mu I$.
\end{proof}

In the computation, $X_k=V|\Lambda|V^\top+\eta_kI$, where $A_k\approx V\Lambda V^\top$ is a floating-point eigendecomposition and $\eta_k\ge\|A_k-V\Lambda V^\top\|_\infty$ is computed rigorously;
then $X_k\mp A_k=V(|\Lambda|\mp\Lambda)V^\top+\bigl(\eta_kI\mp(A_k-V\Lambda V^\top)\bigr)\succeq0$ for any real $V$.
The bound $\bar\rho$ comes from a Collatz--Wielandt quotient, and positive definiteness from an $LDL^\top$ factorization without pivoting.
The affine form of $\Psi$ over $\Theta'$ comes from first-order Taylor (affine) arithmetic with one noise symbol per component of $\theta$.

\begin{theorem}[Strong concavity on $G$]\label{thm:pindyck-concave}
$\nabla^2J(p)\preceq-10^{-3}I$ for every $p\in G$.
\end{theorem}

\begin{proof}
By \cref{lem:pindyck-ranges,lem:pindyck-hessian}, $\nabla^2J(p)=\Psi(\theta(p))$ with $\theta(p)\in\Theta$.
\Cref{lem:pindyck-matrix} with $\mu=10^{-3}$ holds on every leaf of a finite partition of $\Theta$, proved by computer in two ways:
\begin{itemize}
\item the first code: 9 leaves, after bisecting $u_{16}$ once, $u_{15}$ twice, $\beta_{13}$ three times and $\beta_{12}$ twice;
affine models with binary64 coefficients and an analytically derived rounding allowance per operation; $LDL^\top$ in outward-rounded interval arithmetic;
\item the second code: 6 leaves on the hull of both codes' boxes $\Theta$;
affine arithmetic with interval coefficients and one \texttt{nextafter} step per binary64 operation; $LDL^\top$ in exact rational arithmetic.
The second code also passed the first code's 9 leaves.
\end{itemize}
Coverage of $\Theta$ by each partition was checked exactly.
\end{proof}

\paragraph{Proof of \cref{thm:pindyck-bound}.}
Let $p^*$ be the archived price vector with 17 significant digits, and let $g=\nabla J(p^*)$.
The point $p^*$ is the price vector of the exactly feasible point below, and it lies in the interior of $G$ (the smallest slack of $\underline Wp\le\bar\gamma$, computed exactly, is about $6.4376$, and $\min_tp^*_t>9.4$).
For $p\in P\subseteq G$ the segment $[p^*,p]$ lies in the convex set $G$, so Taylor's formula with integral remainder and \cref{thm:pindyck-concave} give
\begin{equation}\label{eq:pindyck-sc}
J(p)\ \le\ J(p^*)+g^\top(p-p^*)-\tfrac{\mu}{2}\|p-p^*\|_2^2\ \le\ J(p^*)+\frac{\|g\|_2^2}{2\mu},\qquad\mu=10^{-3},
\end{equation}
where the second step completes the square.
The second code's interval Newton enclosures at $p^*$ give $J(p^*)\in[J^-,J^+]$ with $J^+-J^-\le\sci{1.4}{-46}$,
\[
J^+=1170.48628543608856208757742506928802541\ldots,
\]
and $\|g\|_2^2\le\sci{3.39889}{-32}$, hence $\|g\|_2^2/(2\mu)\le\sci{1.69945}{-29}$.
By \cref{lem:pindyck-reduction}, every exactly feasible point has objective at least $-J^+-\sci{1.69945}{-29}$, which rounds down to
\[
\Lcert=-1170.486285436088562087577425069306 .
\]
The objective of the exactly feasible point is at most $-J^-$, which rounds up to
\[
\Uprim=-1170.486285436088562087577425069288 .
\]
So $\gapabs\le\sci{1.70}{-29}$ and $\gaprel\le\sci{1.46}{-32}$.
The first code's enclosures ($|\partial J/\partial p_t(p^*)|\le\sci{8.01}{-17}$ for all $t$, enclosure width of $J(p^*)$ about \sci{4.9}{-30}) give the same statement with $\gapabs<\sci{6}{-29}$, since $16\cdot(\sci{8.01}{-17})^2/(2\mu)<\sci{5.14}{-29}$.
Dropping the quadratic term in \eqref{eq:pindyck-sc} and using $0\le p_t\le\bar p_t$ on $G$ gives the weaker bound $-1170.4862854360886163932$ ($\gapabs\le\sci{5.44}{-14}$), which both codes certify.
\hfill$\square$

\paragraph{Proof of \cref{cor:pindyck-unique}.}
The set $P$ is nonempty ($p^*\in P$) and compact (closed, and bounded because $P\subseteq G$), and $J$ is continuous on $G$, so a maximizer $\hat p$ of $J$ on $P$ exists.
By \eqref{eq:pindyck-sc}, $J(\hat p)\ge J(p^*)$ forces $\tfrac{\mu}{2}\|\hat p-p^*\|_2^2\le\|g\|_2\|\hat p-p^*\|_2$, so $\|\hat p-p^*\|_2\le\varrho:=2\|g\|_2/\mu\le\sci{3.7}{-13}$.
The ball $B_\varrho(p^*)$ of radius $\varrho$ around $p^*$ lies in $P$.
Indeed, each row of $\underline W$ has 1-norm at most $\sum_k0.13\cdot0.87^k+\sum_k0.1\cdot0.75^k\le1.4$, so $B_\varrho(p^*)\subseteq G$ by the slack $6.4376$ and $\min_tp^*_t>9.4$.
On $G$, $\iota_t,\phi_t,E_t\in(0,1]$ and $\beta_t\le1.1+0.1\cdot126.15<13.72$, so the recursion of \cref{lem:pindyck-hessian} gives
$\|\nabla s_t\|_1\le a_t$ with $a_0=0$ and $a_t=0.75\,a_{t-1}+0.1+13.72K\sum_{j<t}a_j$; this yields $a_t<1.34$ for all $t\le16$.
With $\|\nabla td_t\|_1\le0.13\sum_k0.87^k<1$ we get $\|\nabla d_t\|_2\le\|\nabla d_t\|_1<2.34$ on $G$, so on $B_\varrho(p^*)$
$d_t\ge d_t(p^*)-2.34\,\varrho>5.4>0$, using $\min_td_t(p^*)=5.437$.
Hence every maximizer $\hat p$ satisfies $\hat p>0$ and $d_t(\hat p)>0$ for all $t$; by continuity it lies in the interior of $P$, so $\nabla J(\hat p)=0$.
For any other $p\in P$, \cref{thm:pindyck-concave} along $[\hat p,p]\subseteq G$ gives $J(p)\le J(\hat p)-\tfrac{\mu}{2}\|p-\hat p\|_2^2<J(\hat p)$, so $\hat p$ is the unique maximizer of $J$ on $P$.

It remains to return to the model.
Let $\hat x$ consist of the prices $\hat p$ and the states given by the recursions of \cref{lem:pindyck-reduction}.
Its states are nonnegative: $s_t,cs_t>0$ because $s_0>0$ and every supply increment is positive, $d_t(\hat p)>0$, $td_t=d_t+s_t>0$, and $R_t\ge192.5$ by \cref{lem:pindyck-ranges}.
So $\hat x$ is exactly feasible, and its objective $-J(\hat p)=-\max_PJ$ is at most the objective of every exactly feasible point (\cref{lem:pindyck-reduction}); hence $\hat x$ is a minimizer.
If $x$ is any minimizer, its prices $p$ lie in $P$ and satisfy $J(p)=J(\hat p)$, so $p=\hat p$, and the recursions then give $x=\hat x$.
\hfill$\square$

\paragraph{Exactly feasible point.}
The prices are $p^*$; every other variable is defined by the recursion of \cref{lem:pindyck-reduction}.
Each implicit $s_t$ is enclosed by a verified interval inclusion: a fixed-point inclusion $T(X)\subseteq X$ in the first code and an interval Newton step in the second.
The point satisfies $\min_td_t=5.437$, $\min_tR_t=350.3$ and $\min_ttd_t=13.99$, so every bound holds.
The decimal vector that approximates the point violates rows by about \sci{8}{-28} and is not itself exactly feasible.

\begin{certbox}[Certificate for \cref{thm:pindyck-bound,cor:pindyck-unique}]
\certitem{Inputs}{OSIL file (SHA-256 \texttt{c7b3b5e4}); LP multipliers; leaf partitions of $\Theta$; 17-digit prices $p^*$.}
\certitem{Computation}{exact LP certificates; matrix test per leaf, \arith{F} and \arith{I}; exact coverage; $J(p^*)$ and $\|g\|_2^2$ in \arith{I}; bound \arith{E}; primal states \arith{I}.}
\certitem{Data reading}{exact or outward; LP and eigenvalue solvers only propose multipliers and vectors.}
\certitem{Implementations}{concavity proved by two separately written codes; the first assumes binary64 round to nearest with an analytically derived rounding allowance per operation, the second \cref{hyp:fp-ieee} with one \texttt{nextafter} step per operation. Displayed bound from the second code's enclosures, a slightly weaker one from the first's; primal point by two codes.}
\certitem{Evidence level}{\evid{stored}.}
\certitem{Replay}{Tier~1; 19\,s and 322\,s for the two concavity proofs.}
\end{certbox}
